
\documentclass[12pt]{article}
\usepackage[utf8]{inputenc}
\usepackage{amsmath, amssymb}
\usepackage{graphicx}
\usepackage{geometry}
\geometry{margin=1in}
\usepackage{caption}
\usepackage{color}
\usepackage{titlesec}
\usepackage{hyperref}

\titleformat{\section}{\large\bfseries}{\thesection}{1em}{}
\titleformat{\subsection}{\normalsize\bfseries}{\thesubsection}{1em}{}

\title{Guión de Clase: Carga Estructural en Apendices Propulsivos}
\author{}
\date{}

\begin{document}

\maketitle

\section*{Objetivo de la clase}
Comprender cómo las fuerzas aerodinámicas inducen momentos estructurales en alas y apéndices propulsivos, desarrollando modelos matemáticos que permitan estimar estos efectos en condiciones reales como el despegue de un albatros.

\section{Introducción (15 minutos)}
\begin{itemize}
    \item Introducción a apéndices propulsivos en la naturaleza (delfines, albatros, samaras).
    \item Analogía entre alas, hélices y estructuras naturales.
    \item Presentación del caso del albatros: ¿Por qué el despegue es crítico?
    \item Visualización: Fig. 5.38 y Fig. 5.42.
\end{itemize}

\section{Desarrollo teórico (30 minutos)}
\subsection{Fuerzas aerodinámicas locales}
\begin{itemize}
    \item Definición de fuerzas: $C_L$, $C_D$, $C_R$, $C_N$, $C_X$.
    \item Ángulo de ataque efectivo: $\alpha_{eff}$.
\end{itemize}

\subsection{Momentos estructurales}
\begin{itemize}
    \item Momento flapwise:
    \[
    M_F = \int_0^R \left( \frac{1}{2} \rho U_{\text{eff}}^2 c \right) C_N z \, dz
    \]
    \item Momento edgewise:
    \[
    M_E = \int_0^R \left( \frac{1}{2} \rho U_{\text{eff}}^2 c \right) C_X z \, dz
    \]
    \item Momento torsional:
    \[
    M_T = \int_0^R \left( \frac{1}{2} \rho U_{\text{eff}}^2 c \right) C_R (x_{EA} - x_{AC}) \, dz
    \]
\end{itemize}

\subsection{Modelos de deformación 2D}
\begin{itemize}
    \item Plunge (movimiento en $y$): Fig. 5.40 izquierda.
    \item Pitch (rotación $\alpha$): Fig. 5.40 derecha.
    \item Inestabilidad de torsión (flutter).
\end{itemize}

\subsection{Torsión y rigidez}
\[
T = JG \frac{d\alpha}{dz}, \quad \alpha = \int_0^{z_1} \frac{T}{JG} dz
\]

\section{Ejemplo práctico: despegue del albatros (40 minutos)}
\begin{itemize}
    \item Datos:
    \begin{itemize}
        \item $b/2 = 1$ m, $c_r = 0.2$ m
        \item $U_\infty = 15$ m/s, $\alpha = 10^\circ$
        \item $C_D = 0.04$, $C_L = 2\pi\alpha$
    \end{itemize}
    \item Distribución elíptica de cuerda:
    \[
    c(z) = c_r \sqrt{1 - 4 \left( \frac{z}{b} \right)^2 }
    \]
    \item Cálculo de $C_R$, $C_N$, $C_X$:
    \[
    C_R = \sqrt{C_L^2 + C_D^2}, \quad C_N = C_L \cos\alpha - C_D \sin\alpha, \quad C_X = C_L \sin\alpha + C_D \cos\alpha
    \]
    \item Momentos:
    \[
    M_F = \frac{c_r b^2}{24} \rho U_\infty^2 C_N \approx 9.839 \text{ Nm}
    \]
    \[
    M_E = \frac{c_r b^2}{24} \rho U_\infty^2 C_X \approx 2.104 \text{ Nm}
    \]
    \[
    M_T = \frac{\pi c_r b}{16} \rho U_\infty^2 C_R (x_{EA} - x_{AC}) \approx 3.554 \text{ Nm}
    \]
\end{itemize}

\section{Cierre y reflexión (5 minutos)}
\begin{itemize}
    \item Conclusiones del caso del albatros.
    \item Importancia de los modelos 2D para analizar estructuras complejas.
    \item Introducción a posibles extensiones: flapping, no linealidades, acoplamiento fluido-estructura.
\end{itemize}

\end{document}
